% =====================================================================
%  Capitolo di esempio: mostra figure, formule chiave e box.
%  Sostituisci il contenuto mantenendo la struttura.
% =====================================================================
\section{Argomento 1}

Breve introduzione all'argomento: a cosa serve e cosa si impara.

\subsection{Una formula chiave}

Le formule importanti si scrivono con \verb|\keyeq{titolo}{formula}{etichetta}|:
vengono numerate, finiscono nell'indice del formulario e nel formulario completo
in fondo alla dispensa.

\keyeq{Legge di Ohm}{V = R\cdot I}{eq:arg1:ohm}

La \eqref{eq:arg1:ohm} lega tensione, corrente e resistenza.

\subsection{Un concetto da ricordare}

\keybox{Resistori in serie}{%
  Sono attraversati dalla stessa corrente; la resistenza equivalente è la somma:
  \[ R_{eq} = R_1 + R_2 + \dots + R_n \]
}{box:arg1:serie}

\subsection{Una figura}

Figure e tabelle restano dove le scrivi (non serve \texttt{[H]}).
Per un'immagine: \verb|\includegraphics[width=0.6\textwidth]{nomefile}|
con il file in \texttt{Dispensa/figure/}.

\begin{figure}
  \centering
  \begin{circuitikz}
    \draw (0,0) to[vsource, l=$V$, n=G] (0,2.5)
          to[R, l=$R_1$] (3,2.5)
          to[R, l=$R_2$] (3,0) -- (0,0);
    \morsetti{G}
  \end{circuitikz}
  \caption{Due resistori in serie.}
  \label{fig:arg1:serie}
\end{figure}
